\chapter{Hypothesis and Objectives}
\label{chapter:hypothesis-objectives}
\epigraph{``If we choose, we can live in a world of comforting illusion''.}{\vspace{10pt}Noam Chomsky}

\newpage
%Example text:

This chapter presents the hypothesis formulated for the thesis. Also, it displays the established objectives with its corresponding deliverables. Last section gives a summary of the scope and limitations defined for this research.

\section{Hypothesis}
\label{section:hypothesis}

%A brief context description of the formulation of the hypothesis can go here, followed by.


The hypothesis for this research effort reads as follows:

%``Our proposed method $P$ in the research can improve by $Q$\footnote{If numbers are involved, usually an explanation of why that specific number was used is given in a footnote \cite{examplereference2}.} the relation in metric 1 and metric 2\footnote{The \cite{examplereference} definition of is used.}.''

%''By systematically exploring and evaluating various modifications to the model's head architecture of a Visual Transformer for Facial Expression
%Recognition (FER), it is possible to surpass the current state-of-the-art performance achieved by the POSTER model on the RAF-DB dataset. Specifically, this study aims to improve the accuracy from POSTER’s reported 92.05\% to at least 93.00\%''

%By systematically exploring and evaluating various modifications to the classification head architecture of a Visual Transformer for Facial Expression Recognition, it is possible to achieve performance improvements using RAF-DB data set, particularly for challenging emotion classes. These improvements are expected to be reflected in increased class-wise precision, recall, and f1-score for under performing categories, while maintaining or improving overall model performance. 

Under a controlled experimental design, modifications to the classification head of a Visual Transformer based Facial Expression Recognition (\gls{FER}) model can produce measurable improvements in class-wise performance metrics (macro‑averaged F1‑score, precision, and recall) on the RAF‑DB dataset, particularly for underrepresented emotion classes, when compared to the original classification head, while maintaining comparable overall accuracy.

\section{Objectives}
\label{section:objectives}
The current research established the following objectives:

\subsection{Main Objective}
\label{section:main-objective}
\label{objective:main}
%Main Objective Text.

To analyze how the design of the classification head in a Visual Transformer based Facial Expression Recognition (\gls{FER}) model influences class discrimination performance in \gls{in-the-wild} conditions, through a controlled experimental study on the RAF‑DB dataset.

\subsection{Specific Objectives}
\label{section:specific-objectives}
% \begin{enumerate}
% 	\item Specific objective 1.
	
% 	\item Specific objective 2.
	
% 	\item Specific objective 3.
	
% \end{enumerate}

\begin{enumerate}

\item Design and implement an enhanced classification head for the POSTER Visual Transformer architecture that incorporates normalization, regularization, and non‑linear transformations, while keeping the backbone unchanged.


\item Evaluate the impact of classification head modifications using a two factor factorial experimental design, comparing models with and without head modifications across three POSTER architecture sizes (Small, Base, and Large).


\item Analyze performance differences between the original and modified classification heads using global metrics (accuracy) and class‑wise oriented metrics  (macro‑averaged F1‑score, precision, and recall).


\item Examine the behavior of the modified classification head on class imbalanced emotion categories in RAF‑DB, with specific attention to underrepresented and difficult expressions.


\item Compare computational efficiency and performance trade‑offs among the evaluated configurations, analyzing the balance between accuracy, macro-avaraged F1-score  and training time using the original and modified classification heads.

\end{enumerate}

\subsection{Deliverables}
%Example text:


The following section provides the deliverables assigned for each objective defined on the previous sections. The table \ref{tab:table_deliverables} shows the corresponding deliverables worked during the development of the thesis proposal and thesis work.



\renewcommand{\arraystretch}{2.5}
\begin{table}[H]
\caption{Deliverables for the Research}
\label{tab:table_deliverables}
\begin{tabular}{@{}llp{10cm}@{}}
\toprule
\textbf{ID} & \textbf{Name} & \textbf{Description} \\ \midrule
COD-01 & Modified POSTER & 
\parbox[t]{10cm}{Modified version of the POSTER ViT with changes in Head architecture} \\
COD-02 & POSTER Training Script & 
\parbox[t]{10cm}{Script for training the modified POSTER ViT with RAFDB database} \\
%COD-03 & Experiment Script & 
%\parbox[t]{10cm}{A Script to automate the data collection, experimental runs and hypothesis testing
%} \\
DOC-01 & Thesis Proposal Document & 
\parbox[t]{10cm}{A proposal thesis work comprising the introduction, theoretical framework, objectives and hypothesis, which has been previously approved for the development of the thesis} \\
DOC-02 & Article Draft & 
\parbox[t]{10cm}{A draft of a scientific article with the research, visual transformer's classification head changes, findings, and conclusions} \\
DOC-03 & Final Thesis Document & 
\parbox[t]{10cm}{A final thesis document with all the previous deliverables' details, including an in-depth academic scientific analysis relevant to this research} \\
PRE-01 & Thesis Proposal Presentation & 
\parbox[t]{10cm}{A presentation encapsulating the scope of the study and a summary of the proposal} \\
PRE-02 & Thesis Defense Presentation & 
\parbox[t]{10cm}{A presentation with all the materials required for the thesis defense} \\ 
\bottomrule
\end{tabular}
\end{table}


The deliverables developed during the research were completed in different stages of the thesis work. During the thesis proposal phase, the deliveries addressed were DOC-01 (Thesis Proposal Document) and PRE-01 ( Thesis Proposal Presentation). During the thesis development and experimentation phase, the deliverables completed include COD‑01 (Modified POSTER implementation), COD‑02 (Training script), COD‑03 (Experiment automation script), DOC‑03 (Final Thesis Document), and PRE‑02 (Thesis Defense Presentation). 


% \textbf{Specific Objective 1:}
% Deliverable for specific objective 1. 

% \textbf{Specific Objective 2:}
% Deliverable for specific objective 2.

% \textbf{Specific Objective 3:}
% Deliverable for specific objective 3.

\section{Scope and Limitations}
\label{section:scope-limitations}


Should include the scope and limitations of the developed research.

% \textbf{Scope and limitation 1:} 

% \textbf{Scope and limitation 2:}
 
% \textbf{Scope and limitation 3:} 

\begin{itemize}
\item This research focuses exclusively on predicting seven discrete class FER emotions classes happy, sad, surprise, fear, disgust, anger, and neutral. It deliberately excludes other FER prediction approaches such as analyzing Facial Action Units (AUs) or the Valence-Arousal (VA), which involve a more nuanced or continuous representation of the emotions \cite{mollahosseini_affectnet_2019, kollias_affect_2021}.

\item This study evaluates the proposed classification head modifications exclusively on the RAF‑DB dataset. While multi‑dataset evaluation can strengthen generalization claims, the use of a single dataset was a design decision to maintain strict experimental control. RAF‑DB was selected due to its widespread adoption in Facial Expression Recognition research. The \gls{in-the-wild} characteristics and its pronounced class imbalance, directly aligns with the objective of analyzing minority-class behavior.

Extending the experimental setup to additional datasets such as AffectNet \cite{mollahosseini_affectnet_2019} or FER2013 \cite{khaireddin_facial_2021} would significantly increase computational cost and experimental complexity. Each additional dataset requires dataset‑specific preprocessing, class remapping, hyperparameter validation, and full retraining across all model sizes and head configurations. Given that this study evaluates six configurations (three model sizes $\times$ two head designs) over $250$ training epochs each, incorporating a second dataset would approximately double the total training time and resource consumption. Average training times in minutes for (Base/Small/Large) with and without changes can be found on table \ref{tab:training_time_experiments}.

\item The project did not include research that aims to analyze emotional states from dynamic sequences for example video-based or video-voice methods \cite{zhang_transformer-based_2022, wang_survey_2024}, only static image-based FER analysis was conducted during the development. While dynamic FER methods may capture temporal emotional transitions more effectively, this project centers on evaluating static images to streamline experimentation and ensure consistency on the experiments.

\item The materials used in this research are limited to licensed and public resources. These include FER data sets, academic publications, source code, programming languages, and libraries, ensuring transparency, reproducibility, and ethical compliance in all aspects of the study.

\end{itemize}